\chapter{Formulário}
\label{cap:formulario}

Este apêndice reúne, em um só lugar, as expressões usadas ao longo do livro. Cada uma
remete ao capítulo em que foi introduzida e discutida.

\section{Métricas de classificação}

Com VP, VN, FP e FN denotando verdadeiros positivos, verdadeiros negativos, falsos
positivos e falsos negativos (\cref{cap:aprendizado}):

\begin{align}
\text{ACC} &= \frac{\text{VP} + \text{VN}}{\text{VP} + \text{VN} + \text{FP} + \text{FN}}
\label{eq:for-acc} \\[0.4em]
\text{REC} &= \frac{\text{VP}}{\text{VP} + \text{FN}}
\label{eq:for-rec} \\[0.4em]
\text{SP} &= \frac{\text{VN}}{\text{VN} + \text{FP}}
\label{eq:for-sp} \\[0.4em]
\text{PR} &= \frac{\text{VP}}{\text{VP} + \text{FP}}
\label{eq:for-pr} \\[0.4em]
F_1 &= 2 \cdot \frac{\text{PR} \cdot \text{REC}}{\text{PR} + \text{REC}}
\label{eq:for-f1} \\[0.4em]
\text{ACC}_b &= \frac{1}{|\mathcal{K}|} \sum_{k \in \mathcal{K}} \text{REC}_k
\label{eq:for-accb}
\end{align}

Para segmentação, com $A$ e $B$ os conjuntos de instantes preditos e anotados:

\begin{equation}
\text{IoU} = \frac{|A \cap B|}{|A \cup B|}
\label{eq:for-iou}
\end{equation}

Para listas ordenadas de $N$ segmentos, com $\rho_i$ a posição da classe correta:

\begin{equation}
\text{top-}k = \frac{1}{N} \sum_{i=1}^{N} \mathbb{1}[\rho_i \le k]
\label{eq:for-topk}
\end{equation}

\section{Erro de reconstrução e limiar}

Erro percentual absoluto médio simétrico, usado como perda e como escore de anomalia
(\cref{cap:dual}):

\begin{equation}
\text{SMAPE} = \frac{100}{n} \sum_{i=1}^{n}
\frac{|\hat{x}_i - x_i|}{(|x_i| + |\hat{x}_i|)/2}
\label{eq:for-smape}
\end{equation}

Limiar sobre o erro de reconstrução, com $\mu$ e $\sigma$ estimados no conjunto de
treinamento e $k = 3$:

\begin{equation}
\tau = \mu + k\,\sigma
\label{eq:for-limiar}
\end{equation}

\section{Detecção de novidade}

Distância de Mahalanobis de um ponto $\vect{x}$ à distribuição de média $\bm{\mu}$ e
covariância $\bm{\Sigma}$ (\cref{cap:mahalanobis}):

\begin{equation}
d_M(\vect{x}) = \sqrt{(\vect{x} - \bm{\mu})^{\top}\,
\bm{\Sigma}^{-1}\,(\vect{x} - \bm{\mu})}
\label{eq:for-mahalanobis}
\end{equation}

Regularização da covariância, necessária quando a dimensão é comparável ao número de
amostras:

\begin{equation}
\bm{\Sigma}_{\text{reg}} = \bm{\Sigma} + \lambda \mat{I}
\label{eq:for-sigmareg}
\end{equation}

Decisão de novidade, com $\tau_{95}$ o percentil 95 das distâncias no conjunto de
referência:

\begin{equation}
\text{novidade}(\vect{x}) =
\begin{cases}
1, & d_M(\vect{x}) > \tau_{95} \\
0, & \text{caso contrário}
\end{cases}
\label{eq:for-decisao}
\end{equation}

Rejeição por ensemble binário (\cref{cap:binarios}), com $s_k$ o escore do classificador
da classe $k$:

\begin{equation}
\text{rejeitar} \iff \max_{k \in \mathcal{K}} s_k < \tau
\label{eq:for-rejeicao}
\end{equation}

Voto ponderado entre classificador latente e OCSVM (\cref{cap:mundoaberto}):

\begin{equation}
v = w_{\text{bin}}\, s_{\text{bin}} + w_{\text{ocsvm}}\, s_{\text{ocsvm}},
\qquad w_{\text{bin}} + w_{\text{ocsvm}} = 1
\label{eq:for-voto}
\end{equation}

\section{Agrupamento}

Objetivo do $K$-médias:

\begin{equation}
\min_{\{C_j\}} \sum_{j=1}^{K} \sum_{\vect{x} \in C_j}
\|\vect{x} - \bm{\mu}_j\|^2
\label{eq:for-kmeans}
\end{equation}

Silhueta de uma amostra, com $a$ a distância média ao próprio grupo e $b$ a distância
média ao grupo vizinho mais próximo:

\begin{equation}
s = \frac{b - a}{\max(a, b)}
\label{eq:for-silhueta}
\end{equation}

Índice de Calinski-Harabasz, com $B$ e $W$ as dispersões entre e dentro dos grupos e $n$
o número de amostras:

\begin{equation}
\text{CH} = \frac{\operatorname{tr}(B)}{\operatorname{tr}(W)} \cdot
\frac{n - K}{K - 1}
\label{eq:for-ch}
\end{equation}

Índice de Davies-Bouldin, com $\sigma_j$ a dispersão do grupo $j$ e $d_{ij}$ a distância
entre centroides:

\begin{equation}
\text{DB} = \frac{1}{K} \sum_{i=1}^{K}
\max_{j \ne i} \frac{\sigma_i + \sigma_j}{d_{ij}}
\label{eq:for-db}
\end{equation}

Limiar adaptativo do ciclo de mundo aberto, com $R = \mu/\sigma$ das distâncias
intragrupo (\cref{cap:mundoaberto}):

\begin{equation}
k(R) =
\begin{cases}
2, & R \ge 4 \\[0.2em]
1 + \dfrac{R - 2}{2}, & 2 \le R < 4 \\[0.6em]
0, & R < 2
\end{cases}
\label{eq:for-kadapt}
\end{equation}

\section{Formação de hidrato}

Temperatura de formação de hidrato pela correlação de Motiee, com $P$ em unidades de
pressão e $\gamma$ a densidade relativa do gás (\cref{cap:closedworld}):

\begin{equation}
T_{\text{hid}} = a_0 + a_1 \log_{10} P + a_2 (\log_{10} P)^2
+ a_3 \gamma + a_4 \gamma \log_{10} P + a_5 \gamma^2
\label{eq:for-motiee}
\end{equation}

Índice de risco analítico, com $T$ a temperatura medida:

\begin{equation}
C_{\text{hid}} = \max\!\left(0,\;
\frac{T_{\text{hid}} - T}{T_{\text{hid}}}\right) \times 100\,\%
\label{eq:for-chid}
\end{equation}

\section{As treze métricas do agente}

Com $\vect{b}$ a janela de linha de base e $\vect{e}$ a janela de evento
(\cref{cap:agente}), para cada sensor:

\begin{align}
\delta &= \bar{e} - \bar{b}
\label{eq:for-delta} \\[0.3em]
\text{pct} &= 100 \cdot \frac{\delta}{|\bar{b}|}
\label{eq:for-pct} \\[0.3em]
z &= \frac{|\delta|}{\max(\sigma_b,\, 10^{-6})}
\label{eq:for-z} \\[0.3em]
\text{noise\_ratio} &= \frac{\sigma_e}{\sigma_b}
\label{eq:for-noise}
\end{align}

Categorizações derivadas:

\begin{table}[htbp]
\centering
\caption{Categorizações das métricas do agente.}
\label{tab:for-categorias}
\small
\begin{tabular}{llp{6.2cm}}
\toprule
\textbf{Métrica} & \textbf{Valor} & \textbf{Condição} \\
\midrule
\multirow{3}{*}{Direção}
& aumento    & $\delta > 0{,}5\,\sigma_b$ \\
& redução    & $\delta < -0{,}5\,\sigma_b$ \\
& estável    & caso contrário \\
\midrule
\multirow{3}{*}{Magnitude}
& pequena    & $z < 1$ \\
& moderada   & $1 \le z < 3$ \\
& grande     & $z \ge 3$ \\
\midrule
\multirow{3}{*}{Taxa}
& súbita     & $|\text{inclinação}| > 0{,}01$ \\
& rápida     & $|\text{inclinação}| > 0{,}002$ \\
& gradual    & caso contrário \\
\bottomrule
\end{tabular}
\end{table}

Métricas restantes: inclinação, comportamento, razão de oscilação, amplitude do evento,
média da linha de base, média do evento, relação entre P-PDG e P-TPT, e sensor mais
afetado.

\section{Intervalo de confiança}

Intervalo de Wilson para uma proporção $\hat{p}$ observada em $n$ ensaios, com
$z_{\alpha/2} = 1{,}96$ para $95\,\%$ (\cref{cap:aprendizado}):

\begin{equation}
\frac{\hat{p} + \dfrac{z^2}{2n} \pm
z \sqrt{\dfrac{\hat{p}(1-\hat{p})}{n} + \dfrac{z^2}{4n^2}}}
{1 + \dfrac{z^2}{n}}
\label{eq:for-wilson}
\end{equation}

\begin{destaque}[Por que Wilson e não a aproximação normal]
Vários resultados do livro envolvem amostras pequenas --- $n = 11$ para a classe 9 no
estudo de novidade --- e proporções próximas de 1, como os $100\,\%$ da classe 2.

A aproximação normal produz, nesses casos, intervalos que ultrapassam os limites $[0,1]$
ou colapsam a largura zero. O intervalo de Wilson permanece válido em ambos os regimes, e
por isso foi usado uniformemente.
\end{destaque}

\section{Normalizações}

Padronização por escore-$z$:

\begin{equation}
x' = \frac{x - \mu}{\sigma}
\label{eq:for-zscore}
\end{equation}

Normalização Min-Max para o intervalo $[0,1]$:

\begin{equation}
x' = \frac{x - \min(x)}{\max(x) - \min(x)}
\label{eq:for-minmax}
\end{equation}
